\documentclass[11pt]{article}
\usepackage[a4paper,margin=23mm]{geometry}
\usepackage{amsmath,amssymb,amsthm,hyperref}
\hypersetup{colorlinks=true,urlcolor=blue,linkcolor=blue}
\newcommand{\R}{\mathbb R}
\newcommand{\Specpar}{\mathcal S}
\newcommand{\id}{\mathrm{id}}
\title{A quadratic operator family with a two-dimensional affine parameter spectrum}
\author{lizaixi01\quad Problem 00000007474}
\date{25 October 2026}
\begin{document}
\maketitle
\vspace{-5mm}

\paragraph{Precise scope.}
The printed conjecture says that the spectrum of a multiparameter quadratic
operator family is a real algebraic curve, followed by an equality involving
genus and pencil rank. Neither language bounds the number of parameters or
specifies a projective quotient. We disprove only the first necessary clause
under the standard \emph{affine parameter-spectrum} convention: the spectrum
consists of parameter vectors at which the operator is not invertible.
We use a real three-parameter family on the genuine Hilbert space $\R$.
We do not compute genus or pencil rank, or refute a projectivized or
otherwise restricted version.

\paragraph{The actual quadratic bounded operators.}
For $x=(s,t,u)\in\R^3$, put
\[
 L(x)=s+t+u,\qquad Q(x)=L(x)^2,\qquad
 T_x=Q(x)\id_{\R}\in\mathcal B_{\R}(\R).
\]
Every $T_x$ is continuous, bounded and selfadjoint. The family depends
continuously on $x$. Its actual coefficient polynomial
$(X_0+X_1+X_2)^2$ is nonzero, homogeneous of degree two, and has total
degree exactly two; in particular this is not a constant family given a
quadratic label. At $(1,0,0)$ the operator is the identity, while at the
origin it is zero.
Multiplication by a real scalar $a$ is invertible exactly when $a\ne0$.
Equivalently, $0$ belongs to its native operator spectrum exactly when
$a=0$. Therefore the actual parameter spectrum is
\[
 \Specpar=\{x:0\in\operatorname{spec}_{\R}(T_x)\}
 =\{(s,t,u):s+t+u=0\}.
\]
This set is nonempty and proper. The argument does not exploit an empty
real locus or an identically singular family.

\paragraph{The genuine vanishing ideal and coordinate ring.}
Let $A=\R[X_0,X_1,X_2]$, $B=\R[Y_0,Y_1]$.
Restrict every polynomial to every actual point of $\Specpar$:
\[
 \rho:A\longrightarrow\R^{\Specpar},\qquad
 \rho(f)(x)=f(x),\qquad I(\Specpar)=\ker\rho.
\]
The coordinate ring is the native quotient $A/I(\Specpar)$.
Define the algebra homomorphism
\[
 \pi:A\longrightarrow B,\quad
 X_0\mapsto Y_0,\ X_1\mapsto Y_1,\ X_2\mapsto-(Y_0+Y_1).
\]
The inclusion $B\to A$, $Y_i\mapsto X_i$ for $i=0,1$, is a right inverse,
so $\pi$ is surjective. Its evaluation identity is
\[
 (\pi f)(a,b)=f(a,b,-a-b).
\]
The map $(a,b)\mapsto(a,b,-a-b)$ covers precisely $\Specpar$.
Consequently $\pi f=0$ implies $f|_{\Specpar}=0$.
Conversely, if $f|_{\Specpar}=0$, then $\pi f$ vanishes at every real
pair. Polynomial evaluation over the infinite integral domain $\R$
separates polynomials, so $\pi f=0$. Hence
\[
 I(\Specpar)=\ker\pi,\qquad A/I(\Specpar)\simeq_{\R}B
\]
by the native first isomorphism theorem for algebras.

\paragraph{Actual Krull dimension.}
A field has Krull dimension zero. The Noetherian polynomial-ring theorem
gives
\[
 \dim\R[Y_0,Y_1]=\dim\R+2=2.
\]
Transport through the proved algebra isomorphism yields
\[
 \operatorname{ringKrullDim}(A/I(\Specpar))=2\ne1.
\]
This is the genuine coordinate-ring Krull dimension, not the vector-space
dimension of a chosen parameterization and not an assigned dimension value.
A real algebraic curve has dimension one, so the actual affine parameter
spectrum above fails the printed curve clause.

\paragraph{Interpretation and formalization boundary.}
The noninvertibility convention for parameter spectra appears explicitly
in Grigorchuk--Yang\,[1]. Yang\,[2] distinguishes the affine preimage
$P(A)$ from its projectivization $p(A)$. Those papers treat linear complex
pencils; here we apply the same pointwise noninvertibility definition to
the genuinely quadratic real family required by the source. They are
definition references, not proofs of our example. The projectivization
of our plane is a different object, so no conclusion about it is claimed.
The usual dimension-one meaning of curve is stated in Stacks,
Definition 33.43.1\,[3]. Lean proves the full operator, actual real-point
restriction, quotient-ring isomorphism and native Krull-dimension chain.
The terminology ``curve'' uses this standard dimension convention; no
general theory of genus or arbitrary schemes is added as a Lean axiom.

\paragraph{Reproduction.}
The accompanying pinned Lean 4 project proves all declarations without
\texttt{sorry}, \texttt{native\_decide} or additional axioms. Run
\texttt{lean --version}, then \texttt{bash verify.sh} for the build,
ordinary exhaustive namespace axiom check, and official fresh kernel replay.
The package records the original commands, logs, source hashes and actual
resource observations. Experimental timings are diagnostic until the
preregistered complete experiment is analyzed.

\begin{thebibliography}{9}\small
\bibitem{gy} R. Grigorchuk and R. Yang,
\emph{Joint spectrum and infinite dihedral group}, arXiv:1605.01547v3 (2017), abstract.
\url{https://arxiv.org/abs/1605.01547}
\bibitem{yang} R. Yang, \emph{Projective spectrum in Banach algebras},
arXiv:0804.0387 (2008), abstract.
\url{https://arxiv.org/abs/0804.0387}
\bibitem{stacks} The Stacks Project, Section 33.43, Definition 33.43.1,
\emph{Curves}, Tag 0A22.
\url{https://stacks.math.columbia.edu/tag/0A22}
\end{thebibliography}
\end{document}
